\documentclass[10pt,a4paper]{amsart}
\usepackage[margin=19mm]{geometry}
\usepackage{amsmath,amssymb,mathtools}
\usepackage[scr=rsfs,cal=euler]{mathalfa}
\usepackage{xfrac}
\usepackage[unicode,hidelinks,bookmarks=false]{hyperref}
\newcommand{\ie}{i.e.\ }
\newcommand{\cf}{cf.\ }
\newcommand{\resp}{respectively\ }
\newcommand{\Subsection}{\subsection}
\providecommand{\nobreakdash}{\nobreak\hbox{-}}
\setlength{\emergencystretch}{2em}
\multlinegap0pt
\usepackage{bm}
\usepackage{amssymb}
\usepackage{float}
\usepackage[shortlabels]{enumitem}
\newtheorem{proposition}{Proposition}
\title{Local and global well-posedness for the nonlinear Schr\"{o}dinger equation with nonhomogeneous boundary conditions}
\author{Engui Fan$^{*, }$, Yuan Li, Xinhan Liu}
\date{Source: arXiv:2606.03064v1 (CC BY 4.0)}
\begin{document}
\maketitle

\noindent\emph{Source and licence.} Local and global well-posedness for the nonlinear Schr\"{o}dinger equation with nonhomogeneous boundary conditions. Version arXiv:2606.03064v1 (CC BY 4.0), distributed by arXiv under the Creative Commons Attribution 4.0 International licence, http://creativecommons.org/licenses/by/4.0/, as recorded in the arXiv licence metadata for this exact version and shown on the abstract page https://arxiv.org/abs/2606.03064v1. Full licence notice: \texttt{licenses/arxiv-2606.03064-CC-BY-4.0.txt}.\par\smallskip

\noindent\textit{Editorial scope.} Counted unit: the article's proposition p33 (source0.tex lines 586--592). The article's complete original proof of it is delivered with it (source0.tex lines 612--667, one \texttt{proof} environment of 56 lines, which the article places away from the statement). No mathematics is written, repaired or re-derived: the statement and the proof are the article's own lines, in the article's own notation, and no retained line is substituted or re-flowed.  Retained uncounted as the source-local context the proof needs: lines 134--135; lines 379--381; lines 409--411; lines 507--508; lines 608--610.  Nothing was omitted from any retained span, so the package carries no omission record.  No retained line contains a citation, so the wrapper carries no bibliography and no bibliography entry.  No retained line contains a figure environment, an image-inclusion command or drawing code, so no image is delivered and the package declares no asset dependency.  Editorial formatting only, in addition: an AMS \texttt{amsart} preamble that restates the article's own declarations and macros used by the retained lines, namely the environments \texttt{proposition} on separate counters, together with the standard packages those lines need.  The retained environments print in this document as Proposition 1 in order of appearance, whereas the article prints the same items with the numbering of its own full document.  The complete native source of the article as distributed in the arXiv e-print for this exact version is bundled unchanged as \texttt{sources/201954/source0.tex}.
\begin{equation}\label{blackmail}			\|f\|_{\mathcal{H}^{s}(\mathbb{R}^{n})}:=\left(\int_{\mathbb{R}^{n}}\left(1+|\xi'|^2+|\xi_{n}|\right)^{s}\left||\xi'|^{2}+\xi_{n}\right|^{\frac{1}{2}}|\widehat{f}(\xi)|^{2}d\xi\right)^{\frac{1}{2}}.
\end{equation}

\begin{equation}\label{thorny}
	w(t)=-i\int_{0}^{t}e^{i(t-\tau)\Delta}f(\tau)d\tau.
\end{equation}

	\begin{equation}\label{cope}	
		\|w\|_{C_{b}(\mathbb{R};\mathcal{H}^{0}(\mathbb{R}^{n}))}\lesssim \|f\|_{L^{\gamma'}(\mathbb{R};L^{\rho'}(\mathbb{R}^{n}))}.
	\end{equation}

\begin{equation}\label{shenjing}			\|f\|_{\mathcal{H}^{s}(\mathbb{R}^{n})}\approx\|f\|_{\mathcal{H}^{s-2}(\mathbb{R}^{n})}+\|\partial_{t}f\|_{\mathcal{H}^{s-2}(\mathbb{R}^{n})}+\|\Delta'f\|_{\mathcal{H}^{s-2}(\mathbb{R}^{n})},\quad s\in \mathbb{R}
\end{equation}

\begin{equation}\label{faction}
	\left\|-i\int_{0}^{t} e^{i(t-\tau)\partial_{x}^{2}}f(\tau)d\tau\right\|_{L_{t}^{2}(\mathbb{R})}\lesssim\|f\|_{L^{1}(\mathbb{R};\dot{H}^{-\frac{1}{2}}(\mathbb{R}))}.
\end{equation}	

\begin{proposition}\label{p33}
	For $s\in\left[0,\frac{3}{2}\right]$, $1\leq \alpha,\beta\leq 2$ satisfing $\frac{1}{\alpha}+\frac{n}{2\beta}=1+\frac{3n-4s}{12}$,  the $w$ given by \eqref{thorny} satisfies
	\begin{equation}\label{taowuji}
		\|w\|_{C_{b}(\mathbb{R};\mathcal{H}^{s}(\mathbb{R}^{n-1}\times(0,T)))}\lesssim\langle T\rangle^{\frac{s}{3}}\|f\|_{L^{\alpha}(0,T;W^{s,\beta}(\mathbb{R}^{n}))}.
	\end{equation}
	In addition, when  $n=2$ and $s\neq\frac{3}{2}$, we require $\left(\alpha,\beta\right)\neq \left(2,\frac{6n}{3n+6-4s}\right)$.
\end{proposition}

\begin{proof}[Proof of Proposition \ref{p33}]
	Let $\boldsymbol{f}(x,t):=\left\{\begin{aligned}&f(x,t), \quad t\in (0,T)\\
		&0, \quad t\in\mathbb{R}\setminus (0,T)\end{aligned}\right.$, i.e., $\boldsymbol{f}$ is the zero extension of $f$. Denote $\boldsymbol{w}(t):=-i \int_{0}^{t}e^{i(t-\tau)\Delta}\boldsymbol{f}(\tau)d\tau$, clearly, $\boldsymbol{w}(t)=w(t)$ for $t\in (0,T)$. In order to prove \eqref{taowuji}, we first consider $s=\frac{3}{2}$, where the only case is $(\alpha,\beta)=(2,2)$. According to \eqref{shenjing}, we have
	{\small\begin{equation}\label{sj} \|\boldsymbol{w}(x',x_{n},t)\|_{\mathcal{H}^{\frac{3}{2}}(\mathbb{R}^{n})}\approx\|\boldsymbol{w}(x',x_{n},t)\|_{\mathcal{H}^{-\frac{1}{2}}(\mathbb{R}^{n})}+\|\partial_{t}\boldsymbol{w}(x',x_{n},t)\|_{\mathcal{H}^{-\frac{1}{2}}(\mathbb{R}^{n})}+\|\Delta'\boldsymbol{w}(x',x_{n},t)\|_{\mathcal{H}^{-\frac{1}{2}}(\mathbb{R}^{n})}.
	\end{equation}}
	It follows from \eqref{cope} that
	\begin{equation}\label{fir} \|\boldsymbol{w}\|_{L^{\infty}(\mathbb{R};\mathcal{H}^{-\frac{1}{2}}(\mathbb{R}^{n}))}\leq \|\boldsymbol{w}\|_{L^{\infty}\left(\mathbb{R};\mathcal{H}^{0}(\mathbb{R}^{n})\right)}\lesssim\|\boldsymbol{f}\|_{L^{1}(\mathbb{R};L^{2}(\mathbb{R}^{n}))}.
	\end{equation}
	For $\partial_{t}\boldsymbol{w}$, we have
	\begin{equation}\label{wf}
		i\partial_{t}\boldsymbol{w}=-\Delta \boldsymbol{w}+\boldsymbol{f}.
	\end{equation}
	Next, we will introduce \eqref{qian} to estimate $\boldsymbol{f}$. By trace theorem,
	\begin{equation} \label{qian}
		\begin{aligned}
			\left \| \|\boldsymbol{f}(x',x_{n},t)\|_{L_{x'}^{2}(\mathbb{R}^{n-1})}  \right\|_{L_{t}^{2}(\mathbb{R})} &= \left \| \|\boldsymbol{f}(x',0+x_{n},t)\|_{L_{x'}^{2}(\mathbb{R}^{n-1})}  \right\|_{L_{t}^{2}(\mathbb{R})}\\&\lesssim   \left\|\|\boldsymbol{f}(\cdot+x_{n}e_{n},t)\|_{H^{\frac{1}{2}+\varepsilon} (\mathbb{R}^{n})}\right\|_{L_{t}^{2}(\mathbb{R})}\leq \|\boldsymbol{f}\|_{L^{2}(\mathbb{R};H^{\frac{3}{2}}(\mathbb{R}^{n}))} .
		\end{aligned}
	\end{equation}
	From \eqref{blackmail}, the embedding   $L^{2}(\mathbb{R}^{n})\hookrightarrow\mathcal{H}^{-\frac{1}{2}}(\mathbb{R}^{n})$ follows, together with \eqref{qian}, we get
	\begin{equation}\label{batter}
		\begin{aligned}	\left\|\|\boldsymbol{f}(\cdot,x_{n},\cdot)\|_{\mathcal{H}^{-\frac{1}{2}}(\mathbb{R}^{n})}\right\|_{L_{x_{n}}^{\infty}(\mathbb{R})}\lesssim\|\boldsymbol{f}\|_{L^{2}(\mathbb{R};H^{\frac{3}{2}}(\mathbb{R}^{n}))}.
		\end{aligned}
	\end{equation}
	
	In order to estimate $\Delta \boldsymbol{w}$ and $\Delta' \boldsymbol{w}$, we consider
	\begin{equation}\label{pw}
		\partial_{j}^{2} \boldsymbol{w} =-i\int_{0}^{t}e^{i(t-\tau)\Delta} \partial_{j}^{2}\boldsymbol{f}(\tau)d\tau,\quad j=1,\cdots,n.
	\end{equation}
	It suffices to give the estimate for one case, for example, for $\partial_{n}^{2}\boldsymbol{w}$. The estimate for other cases is analogous.
	Taking the Fourier transform to the both sides of the equality \eqref{pw} with respect to $x'$ gives
	\begin{equation}\label{star}	\widehat{\partial_{n}^{2}\boldsymbol{w}}^{x'}(\xi',x_{n},t)=-i\int_{0}^{t}e^{-i|\xi'|^{2}(t-\tau)}\left[e^{i(t-\tau)\partial_{x}^{2}}\widehat{\partial_{n}^{2}\boldsymbol{f}}^{x'}(\xi',\cdot,\tau)\right](x_{n}) d\tau.
	\end{equation}
	According to the one-dimensional estimate \eqref{faction}, we get
	\begin{equation*} \left\|\widehat{\partial_{n}^{2}\boldsymbol{w}}^{x'}(\xi',x_{n},\cdot)\right\|_{L^{2}(\mathbb{R})}\lesssim \int_{\mathbb{R}} \left\|\widehat{\partial_{n}^{2}\boldsymbol{f}}^{x'}(\xi',\cdot,\tau)\right\|_{\dot{H}^{-\frac{1}{2}}(\mathbb{R})}d\tau\leq\int_{\mathbb{R}} \left\|\widehat{\boldsymbol{f}}^{x'}(\xi',\cdot,\tau)\right\|_{\dot{H}^{\frac{3}{2}}(\mathbb{R})}d\tau.
	\end{equation*}
	Using  Plancherel's identity and Minkowski's integral inequality, we obtain
	\begin{equation*}		\left\|\partial_{n}^{2}\boldsymbol{w}(\cdot,x_{n},\cdot)\right\|_{L^{2}(\mathbb{R}^{n})}\lesssim\int_{\mathbb{R}}\left\|\left\|\widehat{\boldsymbol{f}}^{x'}(\xi',\cdot,\tau)\right\|_{\dot{H}^{\frac{3}{2}}(\mathbb{R})}\right\|_{L^{2}_{\xi'}(\mathbb{R}^{n-1})}d\tau\leq\int_{\mathbb{R}}\|\boldsymbol{f}(\tau)\|_{H^{\frac{3}{2}}(\mathbb{R}^{n})}d\tau,
	\end{equation*}
	thus, we have
	\begin{equation}\label{x} \left\|\partial_{j}^{2}\boldsymbol{w}(\cdot,x_{n},\cdot)\right\|_{L^{2}(\mathbb{R}^{n})}\leq\|\boldsymbol{f}\|_{L^{1}(\mathbb{R};H^{\frac{3}{2}}(\mathbb{R}^{n}))},\quad j=1,\cdots,n.
	\end{equation}
	From $L^{2}(\mathbb{R}^{n})\hookrightarrow\mathcal{H}^{-\frac{1}{2}}(\mathbb{R}^{n})$ and \eqref{x}, we get
	\begin{equation}\label{del}		\|\Delta'\boldsymbol{w}(\cdot,x_{n},\cdot)\|_{\mathcal{H}^{-\frac{1}{2}}(\mathbb{R}^{n})}+\|\Delta\boldsymbol{w}(\cdot,x_{n},\cdot)\|_{\mathcal{H}^{-\frac{1}{2}}(\mathbb{R}^{n})}\lesssim\|\boldsymbol{f}\|_{L^{1}(\mathbb{R};H^{\frac{3}{2}}(\mathbb{R}^{n}))}.
	\end{equation}
	Combining \eqref{wf}, \eqref{batter} and \eqref{del}, we obtain
	\begin{equation}\label{proficient} \left\|\partial_{t}\boldsymbol{w}\right\|_{L^{\infty}(\mathbb{R};\mathcal{H}^{-\frac{1}{2}}(\mathbb{R}^{n}))}+\left\|\Delta'\boldsymbol{w}\right\|_{L^{\infty}(\mathbb{R};\mathcal{H}^{-\frac{1}{2}}(\mathbb{R}^{n}))}\lesssim  \|\boldsymbol{f}\|_{L^{1}(\mathbb{R};H^{\frac{3}{2}}(\mathbb{R}^{n}))}+\|\boldsymbol{f}\|_{L^{2}(\mathbb{R};H^{\frac{3}{2}}(\mathbb{R}^{n}))}.
	\end{equation}
	By substituting \eqref{fir} and \eqref{proficient} into \eqref{sj}, we have
	\begin{equation}\label{32}
		\begin{aligned}
			\|w\|_{L^{\infty}(\mathbb{R};\mathcal{H}^{\frac{3}{2}}(\mathbb{R}^{n-1}\times(0,T)))}\leq \|\boldsymbol{w}\|_{L^{\infty}(\mathbb{R};\mathcal{H}^{\frac{3}{2}}(\mathbb{R}^{n}))}\lesssim&\,\|\boldsymbol{f}\|_{L^{1}(\mathbb{R};H^{\frac{3}{2}}(\mathbb{R}^{n}))}+\|\boldsymbol{f}\|_{L^{2}(\mathbb{R};H^{\frac{3}{2}}(\mathbb{R}^{n}))}\\
			\lesssim&\,\langle T\rangle^{\frac{1}{2}} \|f\|_{L^{2}(0,T;H^{\frac{3}{2}}(\mathbb{R}^{n}))}.
		\end{aligned}
	\end{equation}
	Therefore, we complete the proof of \eqref{taowuji} for $s=\frac{3}{2}$. The estimate \eqref{taowuji} is then  obtained by  interpolating between \eqref{32} and \eqref{cope}.
\end{proof}

\end{document}
